\chapter{Scalar Path Integrals and Green Functions}
\label{ch:scalar-path-integrals-green-functions}
\ConventionsMixed


\section{Field-Configuration Kernels}

\begin{convention}[Units]
\label{conv:scalar-path-integral-hbar-one}
Set \(\hbar=1\) throughout this chapter.
\end{convention}

\begin{definition}[Regulated equal-time configuration representation]
\label{def:regulated-equal-time-configuration-representation}
Let \(\widehat\phi(t,\vec x)\) be the canonical real scalar field at time
\(t=x^0\).  The notation involving a spatial field configuration is defined
through finite regulators.  Let \(\Lambda\) be a spatial regulator and let
\(E_\Lambda\) be the resulting finite-dimensional real vector space of spatial
field variables.  Its elements are denoted by \(\varphi_\Lambda\), and a
linear functional \(\ell\in E_\Lambda^\vee\) evaluates the regulated field by
\(\ell(\varphi_\Lambda)\).  The regulated configuration Hilbert space is
\[
  \Hilb_\Lambda=L^2(E_\Lambda,\dd_\Lambda\varphi),
\]
where \(\dd_\Lambda\varphi\) is the Lebesgue measure determined by the chosen
coordinates, lattice cell factors, or Fourier-mode normalization.  The
regulated equal-time field operator associated with \(\ell\) acts by
multiplication:
\[
  (\widehat\phi_\Lambda(\ell)\Psi)(\varphi_\Lambda)
  =
  \ell(\varphi_\Lambda)\Psi(\varphi_\Lambda).
\]

A field-configuration generalized vector at time \(t\) and regulator
\(\Lambda\) is the distributional vector \(\ket{\varphi_\Lambda;t}_\Lambda\)
defined by evaluation:
\[
  {}_\Lambda\langle\varphi_\Lambda;t|\Psi\rangle
  =
  \Psi_t(\varphi_\Lambda),
  \qquad
  \Psi\in\mathcal S(E_\Lambda)\subset\Hilb_\Lambda .
\]
Here \(\mathcal S(E_\Lambda)\) is the Schwartz test space after choosing a
linear coordinate system on \(E_\Lambda\).  The resolution of identity is the
weak identity
\[
  \int_{E_\Lambda}\dd_\Lambda\varphi\,
  \ket{\varphi;t}_\Lambda{}_\Lambda\langle\varphi;t|
  =
  1
\]
on \(\mathcal S(E_\Lambda)\).  Accordingly,
\[
  {}_\Lambda\langle\varphi;t_i|\varphi_i;t_i\rangle_\Lambda
  =
  \delta_\Lambda(\varphi-\varphi_i),
\]
where \(\delta_\Lambda\) is the Dirac delta distribution relative to
\(\dd_\Lambda\varphi\).  The continuum symbols
\(\ket{\varphi;t}\), \(\Psi_t[\varphi]\), and \(\delta[\varphi-\varphi_i]\)
mean a compatible family of these regulated objects together with a stated
limiting topology.
\end{definition}

\begin{definition}[Regulated scalar integration conventions]
\label{def:regulated-scalar-integration-conventions}
For a finite-dimensional bosonic scalar regulator, the canonical reference
density is the coordinate density
\[
  D_\Lambda^{\rm ref}\phi=\dd_\Lambda\phi ,
\]
including the lattice cell factors or Fourier-mode normalization chosen with
the regulator.  If the Euclidean quadratic part is
\[
  S_{\rm kin,\Lambda}[\phi]
  =
  \frac12\langle\phi,C_\Lambda^{-1}\phi\rangle_\Lambda
\]
with positive covariance \(C_\Lambda\), the normalized Gaussian reference
measure is
\begin{equation}
  \dd\mu_{C_\Lambda}(\phi)
  =
  Z_{0,\Lambda}^{-1}
  \exp\{-S_{\rm kin,\Lambda}[\phi]\}\,
  D_\Lambda^{\rm ref}\phi .
  \label{eq:gaussian-reference-measure-from-density}
\end{equation}
If \(L_\Lambda[\phi]\) denotes the remaining interaction, the unnormalized full
Euclidean density is
\begin{equation}
  \dd\rho_{\Lambda,S}(\phi)
  =
  \exp\{-L_\Lambda[\phi]\}\,\dd\mu_{C_\Lambda}(\phi)
  =
  Z_{0,\Lambda}^{-1}
  \exp\{-S_{\rm kin,\Lambda}[\phi]-L_\Lambda[\phi]\}\,
  D_\Lambda^{\rm ref}\phi .
  \label{eq:full-density-reference-versus-gaussian}
\end{equation}
Thus \(D_\Lambda^{\rm ref}\phi\), \(\dd\mu_{C_\Lambda}\), and
\(\dd\rho_{\Lambda,S}\) are distinct named objects related by explicit
Radon--Nikodym factors.  In formulas where the older compact notation
\([D\phi]_\Lambda\) appears, it denotes the reference density
\(D_\Lambda^{\rm ref}\phi\) in the finite-dimensional bosonic scalar setting.
Gaussian-measure formulas must use \(\dd\mu_{C_\Lambda}\), and full
source-functional formulas must display the action weight or the density
\(\dd\rho_{\Lambda,S}\).

For fermionic, gauge, BV, Lorentzian oscillatory, dimensional regularization,
and other schemes beyond finite Lebesgue densities, a chapter must specify the
corresponding regulated integration object and its algebraic or analytic
status separately.
\end{definition}

\begin{table}[H]
\centering
\footnotesize
\begin{minipage}{0.94\textwidth}
\textbf{Finite lattice or finite spin system.}
The object is a finite product of site variables, link variables, or spins.
The measure is an ordinary finite product measure, counting measure, compact
Haar measure, or finite-dimensional density.  Thermodynamic and continuum
limits are separate questions.\par\smallskip

\textbf{Finite Fourier-mode cutoff.}
At finite volume and finite mode number, this is a finite-dimensional field
space with a reference density and, after the quadratic part is chosen, a
Gaussian measure.  Sharp momentum cutoffs are useful perturbatively but need
not preserve locality or reflection positivity.\par\smallskip

\textbf{Smooth covariance or Wilsonian cutoff.}
A profile such as
\(C_\Lambda(k)=\chi(k^2/\Lambda^2)/(k^2+m^2)\) defines a regulated covariance.
With a genuine finite regulator it gives a Gaussian measure or density; in
continuum perturbation theory it is a covariance assignment for graphs and
Wilsonian pushforwards.\par\smallskip

\textbf{Heat-kernel, spectral, or point-splitting cutoff.}
These are smoothing kernels, finite-rank projectors, or separated insertion
prescriptions.  They regulate kernels, traces, determinants, or composite
operators; they are not automatically measures on field configurations.\par\smallskip

\textbf{Pauli--Villars auxiliary-field or determinant regulator.}
This is a massive auxiliary-field or determinant-ratio prescription with
specified statistics, coefficients, masses, and symmetry action.  It is not a
measure on the original field space.  If written as a path integral, the
enlarged field space, its density or integration cycle, and the sign or
indefinite-metric structure of the auxiliary sector must be part of the
definition.  The monograph does not use Pauli--Villars as a default
path-integral construction.\par\smallskip

\textbf{Dimensional regularization.}
This is a meromorphic assignment to perturbative graph distributions, with
specified tensor, spinor, polarization, and external-state algebra in
\(D=d-\varepsilon\).  It is not a Borel measure, not a Berezin measure, and
not a measure on a noninteger-dimensional field-configuration space.\par\smallskip

\textbf{BPHZ, momentum subtraction, MS.}
These are subtraction and coordinate schemes imposed on already specified
amplitudes or graph distributions.  They are not by themselves regulators of
a field measure.
\end{minipage}
\caption{Regulator and subtraction classes used in the monograph.  The
catalog names the object that exists before a limit is taken or before a
formal expansion is interpreted, and separates cutoff measures from regulated
covariances, operator-insertion regulators, Pauli--Villars auxiliary
determinants, dimensional regularization, and subtraction prescriptions.}
\label{tab:regulator-integration-status-catalog}
\end{table}

Table~\ref{tab:regulator-integration-status-catalog} controls the shorthand
used later.  A formula written with a lattice spacing, a Fourier cutoff, or a
smooth covariance cutoff may denote an actual finite-regulator integral only
after the finite field space and density have been specified.  A formula
written in dimensional regularization denotes a perturbative meromorphic graph
assignment; it is not an integral over a noninteger-dimensional space of
fields.  A formula described as Pauli--Villars regulated must display the
auxiliary field or determinant data; without that data it is only the name of
a possible subtraction mechanism, not a path-integral measure.  A subtraction
prescription is a further coordinate choice on regulated or meromorphic data,
not a replacement for stating what was regulated.

\begin{definition}[Finite-regulator time-sliced scalar kernel]
\label{def:finite-regulator-time-sliced-scalar-kernel}
At fixed finite-dimensional \(\Lambda\), the field theory is represented by a
quantum-mechanical system with finitely many coordinates.  Choose a time
partition
\[
  t_i=t_0<t_1<\cdots<t_N=t_f.
\]
With boundary fields \(\varphi_0=\varphi_i\) and
\(\varphi_N=\varphi_f\), the finite time-sliced kernel has the form
\[
  K_{\Lambda,N}(\varphi_f,t_f;\varphi_i,t_i)
  =
  \int_{E_\Lambda^{N-1}}
  \prod_{n=1}^{N-1}\dd_\Lambda\varphi_n\,
  \mathcal N_{\Lambda,N}\,
  \exp\bigl[\ii S_{\Lambda,N}(\varphi_0,\ldots,\varphi_N)\bigr].
\]
The normalization factor \(\mathcal N_{\Lambda,N}\), the discrete action
\(S_{\Lambda,N}\), and possible local measure factors are part of the chosen
operator-symbol and time-slicing convention, exactly as in the
finite-dimensional construction of
Chapter~\ref{ch:hamiltonian-evolution-time-sliced-path-integrals}.
\end{definition}

\paragraph{Finite-dimensional Schr\"odinger regulators and Trotter--Kato.}
The finite-dimensional Trotter--Kato and Feynman--Kac theorems recalled in
Chapter~\ref{ch:hamiltonian-evolution-time-sliced-path-integrals} apply to a
spatial regulator only after the regulator has actually reduced the theory to a
finite-dimensional Schr\"odinger problem.  The regulator must have the following
properties.  The configuration space is a finite-dimensional real vector space
\(E_\Lambda\simeq\mathbb R^{M_\Lambda}\), the kinetic energy is a positive
constant-coefficient quadratic form in the conjugate momenta, and the
Hamiltonian on \(L^2(E_\Lambda,\dd_\Lambda\varphi)\) is the Friedrichs
Hamiltonian associated to
\[
  Q_\Lambda[\Psi]
  =
  \frac12
  \int_{E_\Lambda}
  G_\Lambda^{AB}
  \partial_A\overline{\Psi}\,
  \partial_B\Psi\,
  \dd_\Lambda\varphi
  +
  \int_{E_\Lambda}
  U_\Lambda(\varphi)|\Psi(\varphi)|^2\,
  \dd_\Lambda\varphi ,
\]
where \(G_\Lambda^{AB}\) is positive definite and
\(U_\Lambda:E_\Lambda\to\mathbb R\) is locally Kato and bounded from below.
Under these hypotheses the fixed-\(\Lambda\) transition kernels are exactly the
objects governed by the finite-dimensional theorem: the Euclidean transition
kernel is a Wiener or Brownian-bridge expectation on
\(C([0,\tau],E_\Lambda)\) with potential weight
\(\exp[-\int_0^\tau U_\Lambda(\Phi(s))\dd s]\), after the linear change of
variables that puts \(G_\Lambda^{AB}\) in Euclidean form.

This hypothesis is satisfied, for example, by a finite spatial lattice in a
finite box with stable real polynomial potential energy, and by a genuine
finite-mode truncation with a stable real polynomial potential.  It is not
satisfied by a continuum smooth cutoff that leaves infinitely many spatial
degrees of freedom, nor by a formal rule that only modifies the propagator in
a functional integral.  In those cases the phrase ``the Trotter limit for
\(H_\Lambda\)'' has no content until the Hilbert space, domain, closed
quadratic form, and limiting topology have been specified.

Under the finite-dimensional Schr\"odinger-regulator hypotheses just stated,
the Lorentzian
time-sliced kernels of
Definition~\ref{def:finite-regulator-time-sliced-scalar-kernel} have the
corresponding strong operator or distributional limits after smearing boundary
configurations.  In that case
\[
  \lim_{N\to\infty}
  K_{\Lambda,N}(\varphi_f,t_f;\varphi_i,t_i)
  =
  {}_\Lambda\langle\varphi_f;t_i|
  \exp[-\ii H_\Lambda(t_f-t_i)]
  |\varphi_i;t_i\rangle_\Lambda
\]
as a distributional kernel after smearing the boundary configurations.  The
continuum expression
\[
  \langle\varphi_f;t_f|\varphi_i;t_i\rangle
  =
  \int_{\phi(t_i)=\varphi_i}^{\phi(t_f)=\varphi_f}
  [D\phi]\,
  \exp\left(\ii S[\phi]\right)
\]
is a compact notation for this two-stage construction: first a finite regulator
and time slicing, then a stated continuum or perturbative limiting procedure.
If the regulator is instead a Euclidean covariance cutoff, a Euclidean
spacetime lattice action introduced directly, or a purely perturbative cutoff,
the construction requires its own definition: a Gaussian measure with an
interaction, a finite lattice integral, a transfer-matrix or Hamiltonian
construction, a constructive limit, or a formal asymptotic expansion.
At fixed finite-dimensional \(\Lambda\), the same hypotheses identify the
time-sliced kernels with the Trotter products for the self-adjoint Hamiltonian
\(H_\Lambda\).  The strong operator convergence acts on square-integrable
wavefunctions.  Pairing the resulting operator kernel with Schwartz test
functions in the boundary variables gives the distributional kernel displayed
above.  Passing from the regulated family to a continuum symbol is therefore
notation for the specified continuum or perturbative limiting prescription.

\begin{definition}[Compact scalar path-integral notation]
\label{def:compact-scalar-path-integral-notation}
The compact path-integral notation
\[
  Z=\int [D\phi]\,\ee^{\ii S[\phi]}
\]
is used only as an abbreviation for such regulated limits or for asymptotic
expansions derived from them. Its meaning is part of the chosen regularization
and limiting prescription.  A positive Borel measure is one possible
mathematical object for some bosonic Euclidean scalar theories; it is not the
definition of a general QFT path integral and is not implied by the symbol
\([D\phi]\).
\end{definition}

\begin{definition}[Lorentzian oscillatory path-integral notation]
\label{def:lorentzian-oscillatory-path-integral}
Fix a finite spatial regulator \(\Lambda\), a finite time partition, and
coordinates \(u=(u^1,\ldots,u^M)\) on all intermediate field variables.  With
an insertion \(F\), the regulated Lorentzian expression has the form
\[
  I_{\Lambda,N}(F)
  =
  \int_{\mathbb R^M}
  \exp\!\left(\ii S_{\Lambda,N}(u)\right)
  a_{\Lambda,N}(u;F)\,\dd^M u .
\]
This formula is classified as an oscillatory integral or oscillatory
distribution.  One concrete definition is the Abel boundary value
\[
  I_{\Lambda,N}(F)
  =
  \lim_{\epsilon\downarrow0}
  \int_{\mathbb R^M}
  \exp\!\left(\ii S_{\Lambda,N}(u)-\epsilon Q(u)\right)
  a_{\Lambda,N}(u;F)\,\dd^M u ,
\]
whenever the limit exists in the stated distribution topology; here \(Q\) is a
positive definite quadratic form used only to approach the boundary value.
Equivalent finite-dimensional formulations use the Fourier-transform
definition of oscillatory integrals with amplitudes in a specified symbol
class.

For a nondegenerate quadratic phase
\[
  S(u)=\frac12 u^T A u+J^T u,
  \qquad
  A=A^T,\quad \det A\ne0 ,
\]
the normalized Fresnel value is
\[
  \lim_{\epsilon\downarrow0}
  \int_{\mathbb R^M}
  \exp\!\left(
    \frac{\ii}{2}u^TAu+\ii J^Tu-\frac{\epsilon}{2}u^Tu
  \right)\dd^M u
  =
  (2\pi)^{M/2}
  \ee^{\ii\pi\,\operatorname{sgn}(A)/4}
  |\det A|^{-1/2}
  \exp\!\left(-\frac{\ii}{2}J^TA^{-1}J\right),
\]
where
\(\operatorname{sgn}(A)=n_+(A)-n_-(A)\) is the signature.  The phase
\(\ee^{\ii\pi\operatorname{sgn}(A)/4}\) is part of the finite-regulator
definition; in families of quadratic phases its changes are recorded by the
Maslov index.  For a nonlinear phase, perturbation theory is the stationary
phase expansion of the finite-regulator oscillatory distribution around the
chosen critical configurations, with Hessian signatures and contour or
boundary-value data included in the definition.

The continuum Lorentzian symbol
\[
  \int[D\phi]\,\ee^{\ii S[\phi]}
\]
is therefore an oscillatory pseudo-integral: a shorthand for compatible
finite-regulator oscillatory distributions, their Abel or Fourier boundary
values, or their stationary-phase asymptotic expansions.  This classification
is separate from the Euclidean positive-measure construction used in
constructive scalar examples, from Berezin integration for fermionic
variables, and from gauge-fixed, BRST/BV, lattice, or holonomy constructions
for gauge theories.  The relevant analytic frameworks are finite-dimensional
oscillatory-integral theory and stationary phase in the sense of
H\"ormander, the Maslov-index refinement of the Fresnel phase, and the
Albeverio--H\o egh-Krohn theory of infinite-dimensional Fresnel integrals
when such infinite-dimensional oscillatory objects are constructed directly.
\end{definition}

\begin{table}[H]
\centering
\scriptsize
\renewcommand{\arraystretch}{0.96}
\begin{tabular}{@{}>{\raggedright\arraybackslash}p{0.20\textwidth}
  >{\raggedright\arraybackslash}p{0.40\textwidth}
  >{\raggedright\arraybackslash}p{0.32\textwidth}@{}}
\hline
Regime & Rigorous status & Scope of the statement \\
\hline
\(P(\phi)_2\), including \(\phi^4_2\) &
Constructed non-Gaussian Euclidean and Wightman models for bounded-below
polynomial interactions after Wick ordering and the required local
renormalizations. &
This is the basic scalar positive-measure constructive class in two
spacetime dimensions.  Reflection positivity, locality, and the Wightman/OS
structures can be verified in the constructed models. \\
\(\phi^4_3\) &
Constructed non-Gaussian massive scalar models with stable quartic
interaction after mass, coupling, and vacuum-energy counterterms. &
The construction is substantially harder than \(P(\phi)_2\) and is tied to
superrenormalizability in three spacetime dimensions. \\
Selected low-dimensional superrenormalizable scalar--fermion models &
Constructive results exist for specific Yukawa-type and related models after
the ultraviolet cutoff, Wick ordering, fermion determinant/Pfaffian
conventions, and local renormalizations are fixed. &
Each theorem is model-specific.  The label of a Lagrangian family does not
define a uniform constructive theorem; the required hypotheses must be cited
when such a model is used. \\
Two-dimensional Yang--Mills &
Constructed continuum gauge theory for compact gauge group through
heat-kernel/holonomy measures on generalized connections or lattice
projective limits. &
This is a gauge-theory construction, not a scalar Borel measure on ordinary
fields.  The theory has no local propagating gluon degrees of freedom, and its
status should not be extrapolated to \(D=4\). \\
\(\phi^4_D\), \(D\ge5\) &
For the standard reflection-positive lattice/Ising/\(\lambda\phi^4\)
critical scaling problem with positive quartic coupling, the scaling limits
covered by the triviality theorems are Gaussian. &
The theorem concerns that constructive scaling-limit family.  It is a
nonperturbative obstruction to the standard scalar route to an interacting
continuum limit above four dimensions. \\
\(\phi^4_4\) &
For the standard reflection-positive lattice/Ising/\(\lambda\phi^4\)
critical scaling problem, marginal triviality theorems show Gaussian scaling
limits. &
This is sharper than the perturbative Landau-pole argument.  It still leaves
open the logically different question whether another UV-complete
four-dimensional QFT could agree with formal \(\phi^4_4\) perturbation theory
to all orders while differing nonperturbatively. \\
Four-dimensional pure Yang--Mills &
Finite Wilson lattices are rigorous compact-group integrals, and asymptotic
freedom gives the perturbative ultraviolet mechanism.  Continuum existence
with mass gap remains a separate constructive problem. &
The finite-regulator theory is not the desired continuum QFT.  A continuum
limit must produce gauge-invariant Schwinger functions satisfying positivity,
locality, covariance, and the spectral requirements used for reconstruction. \\
Four-dimensional interacting local QFT of scalar/gauge type &
Open constructive problem once one asks for non-Gaussian local fields or
gauge-invariant observables with Wightman/OS positivity, covariance, locality,
and the spectral properties used later in the monograph. &
Perturbative formal power series, finite regulators, and generalized-free
fields do not supply this construction.  The status of each four-dimensional
model must be stated separately: free, perturbative, finite-regulator,
conditional continuum limit, trivial scaling limit, or open problem. \\
\hline
\end{tabular}
\caption{Selected rigorous existence and triviality regimes relevant to the
symbol \([D\phi]\) and to continuum-limit claims.  The table records theorem
classes and open construction problems; it is not a general definition of path
integrals.}
\label{tab:constructive-qft-status-catalog}
\end{table}

\emph{References for Table~\ref{tab:constructive-qft-status-catalog}.}
Standard entry points are E.~Nelson's construction of the two-dimensional
quartic interaction, B.~Simon's
\emph{The \(P(\phi)_2\) Euclidean (Quantum) Field Theory}, and J.~Glimm and
A.~Jaffe, \emph{Quantum Physics: A Functional Integral Point of View}, for
\(P(\phi)_2\) and \(\phi^4_3\) constructive scalar models; B.~Driver,
A.~Sengupta, and T.~L\'evy for rigorous two-dimensional Yang--Mills
constructions; M.~Aizenman and J.~Fr\"ohlich for scalar triviality above four
dimensions; and M.~Aizenman and H.~Duminil-Copin, ``Marginal triviality of the
scaling limits of critical 4D Ising and \(\phi^4_4\) models,'' \emph{Ann. of
Math.} \textbf{194} (2021) 163--235, for the four-dimensional marginal
triviality theorem.
The Glimm--Jaffe functional-integral construction is used here as a reference
for the positive-measure scalar sector.  The organizing framework of this
monograph is broader: a path integral is specified by the mathematical object
assigned by the chosen regulator and limiting prescription, which may be a
Borel measure, an oscillatory distribution, a Berezin functional, a
gauge-fixed or BV integral, a lattice or holonomy construction, or a formal
asymptotic expansion.  A broader catalog of constructive routes and theorem
status appears in Table~\ref{tab:constructive-master-status} of the
constructive volume; the present table records only the examples needed for
the first path-integral distinctions.

Vacuum correlation functions avoid direct use of field-configuration
eigenvectors. For points \(x_a=(x_a^0,\vec x_a)\), the basic Lorentzian
operator distributions are
\[
  \bra{\vac}
  \widehat\phi(x_1)\cdots\widehat\phi(x_n)
  \ket{\vac}.
\]
These are distributions in the spacetime variables, and their analytic
continuations require spectral and growth assumptions.

\section{Complex Time Contours}

\begin{definition}[Complex-time contour with ordered insertions]
\label{def:complex-time-contour-ordered-insertions}
Let \(z^0\) be a complex time coordinate. If a contour \(C\) in the
\(z^0\)-plane is used for the time variable, the corresponding action is
\[
  S_C[\phi]
  =
  \int_C\dd z^0
  \int\dd^{D-1}\vec x\,
  \mathcal L(\phi,\partial_\mu\phi).
\]
Insertions are ordered along the contour. For vacuum correlation functions the
useful contour has insertion times whose imaginary parts are ordered:
\[
  \operatorname{Im}z_1^0
  <
  \operatorname{Im}z_2^0
  <
  \cdots
  <
  \operatorname{Im}z_n^0.
\]
The real parts are free variables. The ordering in imaginary time provides the
analytic separation that will become time ordering after a uniform rotation to
real time.
\end{definition}

Figure~\ref{fig:complex-time-contour-ordering} records the contour-ordering
convention: the contour carries the action, and the imaginary parts of
insertion times impose the analytic ordering used for later time ordering.

\begin{figure}[H]
\centering
\begin{tikzpicture}[>=Stealth,scale=0.92]
  \draw[->,line width=0.7pt] (-3.5,0) -- (4.1,0)
    node[right] {$\operatorname{Re}z^0$};
  \draw[->,line width=0.7pt] (0,-2.45) -- (0,2.85)
    node[above] {$-\operatorname{Im}z^0$};
  \draw[green!50!black,very thick]
    (-3.0,-2.0) .. controls (-2.2,-1.15) and (-1.15,-0.62) ..
    (-0.35,-0.28) .. controls (0.75,0.18) and (1.65,0.42) ..
    (2.65,0.70) .. controls (3.0,0.80) and (3.25,0.92) .. (3.45,1.05);
  \node[green!45!black] at (-2.75,-1.55) {$C$};
  \foreach \x/\y in {-2.35/-1.28,-0.82/-0.36,1.25/0.34,2.75/0.77} {
    \fill[red!70!black] (\x,\y) circle (2.3pt);
  }
  \foreach \y/\lab in {2.10/z_1^0,1.25/z_2^0,0.42/z_3^0,-1.30/z_4^0} {
    \fill[purple!70!black] (0,\y) circle (2.4pt);
    \node[purple!70!black,right=0.1cm] at (0,\y) {$\lab$};
  }
  \node[align=left,anchor=west] at (1.15,-1.70)
    {$S_C=\displaystyle\int_C\dd z^0
      \int\dd^{D-1}\vec x\,\mathcal L$};
\end{tikzpicture}
\caption{Complex-time contour data for ordered vacuum insertions.  The
ordering of imaginary parts determines the analytic domain that rotates to
Lorentzian time ordering.}
\label{fig:complex-time-contour-ordering}
\end{figure}

The figure records only the ordering data. The analytic statement is the
existence of a holomorphic continuation of the relevant correlation function
in the domain where the imaginary parts are ordered.

\begin{hypothesis}[Ordered tube analyticity for scalar vacuum correlators]
\label{hyp:ordered-tube-analyticity-scalar}
For the scalar vacuum \(n\)-point distribution under consideration, assume
there exists a holomorphic function
\[
  \mathcal W_n(z_1,\ldots,z_n),
  \qquad z_a=(z_a^0,\vec x_a),
\]
on the ordered domain
\[
  \operatorname{Im}z_1^0
  <
  \operatorname{Im}z_2^0
  <
  \cdots
  <
  \operatorname{Im}z_n^0
\]
for fixed real spatial points, with distributional boundary value equal to
the Lorentzian Wightman distribution in the corresponding operator order.
The hypothesis includes the growth condition needed for the boundary value to
exist as a tempered or otherwise specified distribution.  In a Wightman theory
this is a consequence of the spectral condition and locality in the usual
tube-domain construction; in a path-integral construction it is an assumption
or theorem of the chosen reconstruction framework.
\end{hypothesis}

\section{Euclidean Functional Expression}

\begin{definition}[Euclidean continuation of a scalar Lagrangian]
\label{def:euclidean-continuation-scalar-lagrangian}
Put the complex time on the Euclidean axis by writing
\[
  z^0=-\ii\tau,
  \qquad
  x_E=(\tau,\vec x)\in\mathbb R^D.
\]
When the analytic continuation of vacuum correlators exists, Euclidean
field-insertion notation is defined inside ordered correlation functions by
\[
  \widehat\phi(z^0,\vec x)\big|_{z^0=-\ii\tau}.
\]
This is a shorthand for the analytic continuation of Lorentzian Hilbert-space
matrix elements; the Euclidean functional variable in \(S_E[\phi]\) remains the
unhatted field \(\phi\).  For Euclidean times ordered as
\[
  \tau_{i_1}>\tau_{i_2}>\cdots>\tau_{i_n},
\]
this notation records the same analytic boundary value that rotates to the
Lorentzian time-ordered product with
\(x_{i_1}^0>x_{i_2}^0>\cdots>x_{i_n}^0\).
For a Lorentzian Lagrangian density
\[
  \mathcal L(\phi,\partial_0\phi,\partial_i\phi),
\]
define the Euclidean Lagrangian density by
\[
  \mathcal L_E(\phi,\partial_\tau\phi,\partial_i\phi)
  =
  -\mathcal L(\phi,\ii\partial_\tau\phi,\partial_i\phi),
\]
when this analytic substitution is meaningful for the regulated theory. The
Euclidean action is
\[
S_E[\phi]
  =
  \int\dd^Dx_E\,\mathcal L_E(\phi,\partial_E\phi).
\]
\end{definition}

\paragraph{Euclidean scalar action.}
For the real scalar field with
\[
  \mathcal L
  =
  -\frac12\eta^{\mu\nu}\partial_\mu\phi\,\partial_\nu\phi
  -
  V(\phi),
\]
this gives
\[
  \mathcal L_E
  =
  \frac12(\partial_\tau\phi)^2
  +
  \frac12\sum_{i=1}^{D-1}(\partial_i\phi)^2
  +
  V(\phi).
\]
With mostly-plus Lorentzian metric,
\[
  -\frac12\eta^{\mu\nu}\partial_\mu\phi\,\partial_\nu\phi
  =
  \frac12(\partial_0\phi)^2
  -
  \frac12\sum_{i=1}^{D-1}(\partial_i\phi)^2 .
\]
Substitution of \(\partial_0=\ii\partial_\tau\) gives
\[
  \mathcal L(\phi,\ii\partial_\tau\phi,\partial_i\phi)
  =
  -\frac12(\partial_\tau\phi)^2
  -
  \frac12\sum_{i=1}^{D-1}(\partial_i\phi)^2
  -
  V(\phi).
\]
The definition \(\mathcal L_E=-\mathcal L|_{\partial_0=\ii\partial_\tau}\)
then gives the displayed positive kinetic terms.

\begin{definition}[Regulated Euclidean scalar Green function]
\label{def:regulated-euclidean-scalar-green-function}
With a regulator in place, the normalized Euclidean \(n\)-point function is
\[
  \langle\phi(x_{E,1})\cdots\phi(x_{E,n})\rangle_{E,\Lambda}
  =
  \frac{1}{Z_{E,\Lambda}}
  \int D_\Lambda^{\rm ref}\phi\,
  \ee^{-S_{E,\Lambda}[\phi]}\,
  \phi(x_{E,1})\cdots\phi(x_{E,n}),
\]
where
\[
  Z_{E,\Lambda}
  =
  \int D_\Lambda^{\rm ref}\phi\,
  \ee^{-S_{E,\Lambda}[\phi]}.
\]
Here the Euclidean action \(S_{E,\Lambda}\) includes its quadratic part and
\(D_\Lambda^{\rm ref}\phi\) is the reference density of
Definition~\ref{def:regulated-scalar-integration-conventions}.  If
\[
  S_{E,\Lambda}=S_{\rm kin,\Lambda}+L_\Lambda ,
\]
the same normalized correlation function may equivalently be written with the
Gaussian measure \(\dd\mu_{C_\Lambda}\) and interaction weight
\(\ee^{-L_\Lambda[\phi]}\).  The choice of notation therefore records which
factor carries the quadratic kinetic form.
The continuum Euclidean Green function is the limit of these distributions
when the limit exists in the chosen topology and with the chosen
renormalization conditions.
\end{definition}

\section{Uniform Wick Rotation}

\begin{definition}[Uniform Wick-rotation path for ordered times]
\label{def:uniform-wick-rotation-path}
Fix real numbers
\[
  t_1>t_2>\cdots>t_n.
\]
For \(0<\alpha\le\pi/2\), set
\[
  z_a^0(\alpha)=\ee^{-\ii\alpha}t_a.
\]
At \(\alpha=\pi/2\), the insertion lies on the Euclidean axis:
\[
  z_a^0=-\ii t_a.
\]
As \(\alpha\to0^+\), the insertion approaches the Lorentzian time \(t_a\).
Since
\[
  \operatorname{Im}z_a^0(\alpha)=-t_a\sin\alpha,
\]
the ordering
\[
  \operatorname{Im}z_1^0(\alpha)
  <
  \cdots
  <
  \operatorname{Im}z_n^0(\alpha)
\]
is maintained throughout the rotation.
\end{definition}

Figure~\ref{fig:uniform-wick-rotation-insertion-times} displays the uniform
rotation of all insertion times and the preserved imaginary-time ordering along
the rotation path.

\begin{figure}[H]
\centering
\begin{tikzpicture}[>=Stealth,scale=0.9]
  \draw[->,line width=0.7pt] (-1.0,0) -- (5.05,0)
    node[right] {$\operatorname{Re}z^0$};
  \draw[->,line width=0.7pt] (0,-2.35) -- (0,2.65)
    node[above] {$-\operatorname{Im}z^0$};
  \foreach \y/\lab in {2.0/z_1^0,1.15/z_2^0,0.34/z_3^0,-1.20/z_4^0} {
    \fill[purple!70!black] (0,\y) circle (2.5pt);
    \node[purple!70!black,left=0.12cm] at (0,\y) {$\lab$};
  }
  \foreach \starty/\endx/\endy in {2.0/3.3/0,1.15/2.15/0,0.34/0.85/0,-1.20/-0.45/0} {
    \draw[orange!85!black,very thick,->]
      (0.13,\starty) .. controls (0.65,\starty) and (\endx,\endy+0.55) ..
      (\endx,\endy+0.05);
  }
  \foreach \x in {3.3,2.15,0.85,-0.45} {
    \fill[red!70!black] (\x,0) circle (2.3pt);
  }
  \node[align=left,blue!70!black] at (3.25,1.48)
    {$z_a^0(\alpha)=\ee^{-\ii\alpha}t_a$\\
     \(\alpha:\pi/2\to0^+\)};
\end{tikzpicture}
\caption{Uniform Wick rotation of ordered insertion times.  The deformation
\(z_a^0(\alpha)=\ee^{-\ii\alpha}t_a\) preserves the imaginary-time ordering
until the Lorentzian boundary is reached.}
\label{fig:uniform-wick-rotation-insertion-times}
\end{figure}

\paragraph{Euclidean boundary value and the time-ordered function.}
Under Hypothesis~\ref{hyp:ordered-tube-analyticity-scalar}, the Euclidean
correlation function continued along this common rotation becomes the
Lorentzian time-ordered correlator with \(t_1>\cdots>t_n\):
\[
  \langle\phi(x_{E,1})\cdots\phi(x_{E,n})\rangle_E
  \longrightarrow
  \bra{\vac}
  \operatorname{T}\{\widehat\phi(x_1)\cdots\widehat\phi(x_n)\}
  \ket{\vac}.
\]
Here \(\operatorname{T}\) denotes time ordering of the operator-valued
distributions.

By Definition~\ref{def:uniform-wick-rotation-path}, the points
\(z_a^0(\alpha)=\ee^{-\ii\alpha}t_a\) stay inside the ordered domain of
Hypothesis~\ref{hyp:ordered-tube-analyticity-scalar} for every
\(0<\alpha\le\pi/2\).  At \(\alpha=\pi/2\), the points lie on the Euclidean
axis and give the Euclidean boundary value.  As \(\alpha\to0^+\), the ordered
tube boundary value is the Lorentzian Wightman distribution in the operator
order \(1,\ldots,n\).  Since the real times already obey
\(t_1>\cdots>t_n\), that Wightman order is exactly the time-ordered product.

\section{The Free Euclidean Gaussian Measure}

\begin{definition}[Finite-regulator Euclidean Gaussian datum]
\label{def:finite-regulator-euclidean-gaussian-datum}
At finite Euclidean regulator, the free scalar path integral is an ordinary
Gaussian integral.  Let \(F_\Lambda\) be the finite-dimensional real vector
space of regulated Euclidean fields, let \(\dd_\Lambda\phi\) be the regulator
Lebesgue measure, and let
\[
  A_\Lambda:F_\Lambda\longrightarrow F_\Lambda^\vee
\]
be a symmetric positive-definite operator.  The pairing between
\(F_\Lambda^\vee\) and \(F_\Lambda\) is denoted by
\(\langle J,\phi\rangle_\Lambda\).  The regulated free action is
\[
  S_{E,0,\Lambda}[\phi]
  =
  \frac12\langle A_\Lambda\phi,\phi\rangle_\Lambda .
\]
The corresponding normalized Gaussian expectation is
\[
  \langle F(\phi)\rangle_{0,\Lambda}
  =
  \frac{1}{Z_{0,\Lambda}}
  \int_{F_\Lambda}\dd_\Lambda\phi\,
  \ee^{-\frac12\langle A_\Lambda\phi,\phi\rangle_\Lambda}
  F(\phi),
\]
with
\[
  Z_{0,\Lambda}
  =
  \int_{F_\Lambda}\dd_\Lambda\phi\,
  \ee^{-\frac12\langle A_\Lambda\phi,\phi\rangle_\Lambda}.
\]
The covariance is the inverse map
\[
  C_\Lambda=A_\Lambda^{-1}:F_\Lambda^\vee\to F_\Lambda .
\]
\end{definition}

\paragraph{Finite-dimensional Gaussian source functional.}
For sources \(J\in F_\Lambda^\vee\), completing the square gives
\[
  -\frac12\langle A_\Lambda\phi,\phi\rangle_\Lambda
  +\langle J,\phi\rangle_\Lambda
  =
  -\frac12
  \langle A_\Lambda(\phi-C_\Lambda J),\phi-C_\Lambda J\rangle_\Lambda
  +
  \frac12\langle J,C_\Lambda J\rangle_\Lambda .
\]
Therefore
\[
  \frac{Z_{0,\Lambda}[J]}{Z_{0,\Lambda}[0]}
  =
  \exp\left(\frac12\langle J,C_\Lambda J\rangle_\Lambda\right),
  \qquad
  Z_{0,\Lambda}[J]
  =
  \int_{F_\Lambda}\dd_\Lambda\phi\,
  \ee^{-S_{E,0,\Lambda}[\phi]+\langle J,\phi\rangle_\Lambda}.
\]
Differentiating at \(J=0\) gives, for \(f,g\in F_\Lambda^\vee\),
\[
  \bigl\langle
  \langle f,\phi\rangle_\Lambda
  \langle g,\phi\rangle_\Lambda
  \bigr\rangle_{0,\Lambda}
  =
  \langle f,C_\Lambda g\rangle_\Lambda ,
\]
and all higher moments are sums over pairings with this covariance.  This is
Wick's theorem as a finite-dimensional Gaussian identity.

The completion-of-the-square identity follows from
\(A_\Lambda C_\Lambda=\id_{F_\Lambda^\vee}\).  Translation invariance of the
Lebesgue density gives
\[
  \int_{F_\Lambda}\dd_\Lambda\phi\,
  \ee^{-\frac12\langle A_\Lambda\phi,\phi\rangle_\Lambda+\langle J,\phi\rangle_\Lambda}
  =
  \ee^{\frac12\langle J,C_\Lambda J\rangle_\Lambda}
  \int_{F_\Lambda}\dd_\Lambda\phi\,
  \ee^{-\frac12\langle A_\Lambda\phi,\phi\rangle_\Lambda},
\]
after the shift \(\phi\mapsto\phi+C_\Lambda J\).  Dividing by
\(Z_{0,\Lambda}[0]\) gives the source functional.  Applying
\(\partial_s\partial_t\) to the identity with \(J=sf+tg\) and setting
\(s=t=0\) gives the two-point formula.  Repeated differentiation gives the
sum over pairings because the exponent is quadratic in \(J\); odd derivatives
vanish at \(J=0\), and even derivatives contract derivatives in pairs.
This finite-dimensional calculation is the regulator-level meaning of the free
Euclidean Gaussian path integral.  The nontrivial continuum questions are the
choice of regulator, the topology in which \(C_\Lambda\) converges, and the
existence of limiting distributions after the test functions have been fixed.

\begin{definition}[Continuum free scalar Euclidean covariance]
\label{def:continuum-free-scalar-euclidean-covariance}
For the continuum free scalar field of mass \(m\ge0\),
\[
  V(\phi)=\frac12m^2\phi^2,
\]
and the Euclidean action is
\[
  S_{E,0}[\phi]
  =
  \frac12
  \int\dd^Dx_E\,
  \phi(x_E)(-\partial_E^2+m^2)\phi(x_E),
\]
after integration by parts with the chosen boundary conditions. The operator
\[
  A_E=-\partial_E^2+m^2
\]
is the continuum counterpart of \(A_\Lambda\).  If \(m>0\), or if boundary
conditions or an infrared regulator remove the constant zero mode when
\(m=0\), the inverse covariance is represented by the Green function
\(\Delta_E\) defined by
\[
  A_E\Delta_E(x_E-y_E)=\delta^{(D)}(x_E-y_E).
\]
For test functions \(f,g\), the continuum covariance is the distributional
limit
\[
  \lim_{\Lambda}
  \langle f_\Lambda,C_\Lambda g_\Lambda\rangle_\Lambda
  =
  \int\dd^Dx_E\,\dd^Dy_E\,
  f(x_E)\Delta_E(x_E-y_E)g(y_E),
\]
where \(f_\Lambda,g_\Lambda\) are the regulated images of \(f,g\).  On
\(\mathbb R^D\), translation invariance gives the momentum-space kernel
\[
  \Delta_E(x_E-y_E)
  =
  \int\frac{\dd^Dk_E}{(2\pi)^D}
  \frac{\ee^{\ii k_E\cdot(x_E-y_E)}}{k_E^2+m^2},
\]
where
\[
  k_E^2=(k_\tau)^2+\vec k^2,
  \qquad
  k_E\cdot x_E=k_\tau\tau+\vec k\cdot\vec x.
\]
In point notation, after the distribution has been paired with test functions,
the free two-point function is written
\[
  \langle\phi(x_E)\phi(y_E)\rangle_E=\Delta_E(x_E-y_E)
\]
for the free theory.  This equation is precisely the continuum form of the
finite-dimensional covariance identity above.
\end{definition}

\section{Spectral Zeta Determinants}
\label{sec:spectral-zeta-determinants}

The finite-dimensional Gaussian normalization contains the determinant of the
quadratic operator:
\[
  \int_{\mathbb R^N}\dd^N\phi\,
  \exp\left[-\frac12\phi^iA_{ij}\phi^j\right]
  =
  (2\pi)^{N/2}(\det A)^{-1/2}
\]
for a positive-definite symmetric matrix \(A\).  In field theory the analogous
quantity is a determinant of a differential operator and requires its own
regularized definition.  A spectral zeta determinant is one such definition
for operators with discrete positive spectrum.  It regulates the determinant;
it does not by itself choose the continuum theory or the subtraction
coordinates of an interacting path integral.

\begin{definition}[Spectral zeta determinant]
\label{def:spectral-zeta-determinant}
Let \(A\) be a positive self-adjoint elliptic operator of order \(r>0\) on a
compact \(d\)-dimensional Euclidean manifold, with boundary condition chosen so
that \(A\) has compact resolvent.  Let
\[
  0<\lambda_1\le\lambda_2\le\cdots
\]
be the positive eigenvalues of \(A\), repeated with multiplicity.  Zero modes,
if present, are omitted from this list and must be integrated or treated as
collective coordinates separately.  For \(\operatorname{Re}s>d/r\), set
\begin{equation}
  \zeta_A(s)=\sum_j \lambda_j^{-s}.
  \label{eq:spectral-zeta-definition}
\end{equation}
Equivalently, in the same half-plane,
\begin{equation}
  \zeta_A(s)
  =
  \frac{1}{\Gamma(s)}
  \int_0^\infty \dd t\,t^{s-1}
  \operatorname{Tr}\ee^{-tA}.
  \label{eq:spectral-zeta-heat-kernel-mellin}
\end{equation}
Assume that the heat trace has the usual small-\(t\) asymptotic expansion and
that the meromorphic continuation of \(\zeta_A(s)\) is regular at \(s=0\).
For a reference scale \(\mu\) of mass dimension \(1\), define
\begin{equation}
  \log\det\nolimits_{\zeta}\left(\frac{A}{\mu^r}\right)
  =
  -
  \left.\frac{\dd}{\dd s}\right|_{s=0}
  \sum_j\left(\frac{\lambda_j}{\mu^r}\right)^{-s}
  =
  -\zeta_A'(0)-\zeta_A(0)\log\mu^r .
  \label{eq:zeta-determinant-definition}
\end{equation}
\end{definition}

\paragraph{Scale dependence of a zeta determinant.}
In the setting of Definition~\ref{def:spectral-zeta-determinant}, suppose
\[
  \operatorname{Tr}\ee^{-tA}
  \sim
  \sum_{j\ge0} a_j(A)\,t^{(j-d)/r}
  \qquad (t\downarrow0)
\]
with no logarithmic term at order \(t^0\), and suppose zero modes have been
removed.  Then
\[
  \zeta_A(0)=a_d(A).
\]
Consequently a change \(\mu\mapsto \ee^\sigma\mu\) changes the determinant by
\[
  \log\det\nolimits_\zeta\!\left(\frac{A}{(\ee^\sigma\mu)^r}\right)
  -
  \log\det\nolimits_\zeta\!\left(\frac{A}{\mu^r}\right)
  =
  -r\sigma\,a_d(A).
\]
Thus the \(\mu\)-dependence of the one-loop determinant is controlled by the
local heat-kernel coefficient \(a_d(A)\).
The Mellin representation
\[
  \zeta_A(s)=\frac1{\Gamma(s)}
  \int_0^\infty \dd t\,t^{s-1}\operatorname{Tr}\ee^{-tA}
\]
continues meromorphically by subtracting the small-\(t\) heat-kernel
asymptotic series.  The coefficient \(a_d(A)\) multiplies \(t^0\), hence
contributes \(a_d(A)/(s\Gamma(s))\) to \(\zeta_A(s)\) near \(s=0\).  Since
\(\Gamma(s)^{-1}=s+O(s^2)\), this contribution has value \(a_d(A)\) at
\(s=0\).  The remaining subtracted terms are either regular and vanish at
\(s=0\) after multiplication by \(\Gamma(s)^{-1}\), or have poles away from
\(s=0\) under the stated no-log hypothesis.  This gives
\(\zeta_A(0)=a_d(A)\).  Substituting
\((\ee^\sigma\mu)^r=\ee^{r\sigma}\mu^r\) into the defining formula gives
\[
  \log\det\nolimits_{\zeta}\left(\frac{A}{(\ee^\sigma\mu)^r}\right)
  =
  -\zeta_A'(0)-\zeta_A(0)(\log\mu^r+r\sigma),
\]
and subtracting
\[
  \log\det\nolimits_{\zeta}\left(\frac{A}{\mu^r}\right)
  =
  -\zeta_A'(0)-\zeta_A(0)\log\mu^r
\]
leaves \(-r\sigma\zeta_A(0)=-r\sigma a_d(A)\).

This is the determinant counterpart of the regulator/subtraction distinction
recorded in Table~\ref{tab:regulator-integration-status-catalog}: changing
the determinant scale changes a one-loop effective action by a local
heat-kernel coefficient, not by a new nonlocal observable.

\paragraph{One-loop determinant from a quadratic fluctuation.}
Let \(\phi=\phi_{\rm cl}+\eta\) be a background split for a bosonic Euclidean
path integral, and suppose that the quadratic fluctuation operator
\[
  A_{\phi_{\rm cl}}
  =
  S_E''[\phi_{\rm cl}]
\]
is positive and satisfies the hypotheses of
Definition~\ref{def:spectral-zeta-determinant}.  With the Gaussian
normalization chosen by spectral zeta regularization, the one-loop
contribution of the nonzero modes to the effective action is
\begin{equation}
  \Gamma_{1}[\phi_{\rm cl}]
  =
  \frac12
  \log\det\nolimits_{\zeta}\left(\frac{A_{\phi_{\rm cl}}}{\mu^r}\right).
  \label{eq:zeta-one-loop-effective-action}
\end{equation}
For fermionic Gaussian variables the corresponding Pfaffian or determinant
appears with the opposite statistics sign and with the spinor and reality
structure stated as part of the regulator.
At finite regulator,
\[
  S_E[\phi_{\rm cl}+\eta]
  =
  S_E[\phi_{\rm cl}]
  +
  \frac12\langle\eta,A_{\phi_{\rm cl},\Lambda}\eta\rangle_\Lambda
  +
  O(\eta^3)
\]
when \(\phi_{\rm cl}\) obeys the regulated equation of motion.  The Gaussian
integral over the nonzero fluctuation modes gives
\[
  Z_{\Lambda}^{\rm 1-loop}[\phi_{\rm cl}]
  =
  \exp[-S_E[\phi_{\rm cl}]]
  (2\pi)^{N_\Lambda/2}
  \bigl(\det A_{\phi_{\rm cl},\Lambda}\bigr)^{-1/2}
\]
up to the finite-dimensional zero-mode and collective-coordinate factors that
were excluded.  The factor \((2\pi)^{N_\Lambda/2}\) is independent of the
background in a fixed bosonic coordinate density and is part of the vacuum
normalization.  Replacing the divergent determinant by
\(\det_\zeta(A_{\phi_{\rm cl}}/\mu^r)\) gives the logarithm
\eqref{eq:zeta-one-loop-effective-action}.  The argument uses only the
quadratic Gaussian identity; interactions beyond quadratic order are
additional vertices in perturbation theory.

\begin{example}[Thermal harmonic oscillator determinant]
\label{ex:zeta-thermal-harmonic-oscillator}
Let \(S^1_\beta=\mathbb R/\beta\mathbb Z\) and
\[
  A_\omega=-\frac{\dd^2}{\dd\tau^2}+\omega^2,
  \qquad
  \omega>0,
\]
acting on periodic functions.  The eigenvalues are
\[
  \lambda_n=\left(\frac{2\pi n}{\beta}\right)^2+\omega^2,
  \qquad n\in\mathbb Z.
\]
For this one-dimensional second-order operator \(\zeta_{A_\omega}(0)=0\), so
the determinant is independent of \(\mu\).  Differentiate the logarithm before
integrating it:
\[
  \frac{\dd}{\dd\omega}
  \log\det\nolimits_{\zeta}A_\omega
  =
  2\omega
  \sum_{n\in\mathbb Z}
  \frac{1}{(2\pi n/\beta)^2+\omega^2}.
\]
The partial-fraction identity
\[
  \sum_{n\in\mathbb Z}\frac{1}{n^2+a^2}
  =
  \frac{\pi}{a}\coth(\pi a),
  \qquad a>0,
\]
with \(a=\beta\omega/(2\pi)\) gives
\[
  \frac{\dd}{\dd\omega}
  \log\det\nolimits_{\zeta}A_\omega
  =
  \beta\coth\frac{\beta\omega}{2}.
\]
Therefore
\[
  \log\det\nolimits_{\zeta}A_\omega
  =
  2\log\left(2\sinh\frac{\beta\omega}{2}\right)+C .
\]
The constant is fixed by the small-\(\omega\) behavior:
\(\det_\zeta A_\omega\sim \omega^2\det_\zeta'(-\dd^2/\dd\tau^2)\) and
\(\det_\zeta'(-\dd^2/\dd\tau^2)=\beta^2\).  Hence
\begin{equation}
  \det\nolimits_{\zeta}A_\omega
  =
  4\sinh^2\frac{\beta\omega}{2},
  \qquad
  Z_{\rm osc}(\beta,\omega)
  =
  \left(\det\nolimits_{\zeta}A_\omega\right)^{-1/2}
  =
  \frac{1}{2\sinh(\beta\omega/2)} .
  \label{eq:zeta-oscillator-determinant}
\end{equation}
The last expression is the canonical thermal partition function of the
oscillator,
\[
  \operatorname{Tr}\ee^{-\beta\omega(a^\dagger a+1/2)}
  =
  \frac{\ee^{-\beta\omega/2}}{1-\ee^{-\beta\omega}}.
\]
\end{example}

\begin{example}[Circle Casimir finite part]
\label{ex:zeta-circle-casimir-energy}
For a massless real scalar on a spatial circle of circumference \(L\), remove
the constant zero mode and write the oscillator frequencies as
\[
  \omega_n=\frac{2\pi |n|}{L},\qquad n\in\mathbb Z\setminus\{0\}.
\]
The zeta-regularized finite part of the vacuum energy is
\begin{equation}
  E_0^{\zeta}(L)
  =
  \frac12\sum_{n\ne0}\omega_n
  =
  \frac{2\pi}{L}\zeta_{\rm R}(-1)
  =
  -\frac{\pi}{6L},
  \label{eq:zeta-circle-casimir}
\end{equation}
where \(\zeta_{\rm R}(-1)=-1/12\) is the analytic continuation of the Riemann
zeta function.  The zero mode and any additive local vacuum-energy
normalization are separate data.  In two-dimensional CFT language the result
is the \(c=1\) case of \(E_0=-\pi c/(6L)\) on the cylinder.
\end{example}

\section{False-Vacuum Decay as a Contour Problem}
\label{sec:false-vacuum-decay-contour-problem}

The Euclidean word ``bounce'' is dangerous when it is separated from the
state and contour data that make a decay rate.  A stationary point of
\(S_E\) is not by itself a physical width.  The construction has four
distinct layers: a metastable state or branch, a complex integration cycle, a
negative fluctuation direction, and a real-time interpretation.  The standard
semiclassical formulae of Coleman and Callan--Coleman
\cite{Coleman1977FalseVacuum,CallanColeman1977Fate} are used below as
historical anchors; the logical content here is the finite-regulator
calculus, not the authority of the name.

\begin{definition}[Finite-volume false-vacuum datum]
\label{def:finite-volume-false-vacuum-datum}
At fixed finite spatial volume and regulator let \(H_\Lambda\) be a
self-adjoint Hamiltonian on \(\Hilb_\Lambda\).  Since
\(\operatorname{spec}(H_\Lambda)\subset\mathbb R\), a false vacuum is not an
exact decaying eigenvector.  A finite-volume false-vacuum datum consists of:
\begin{enumerate}
\item a normalized vector or density matrix \(\Psi_F\) localized in the basin
      of a local minimum, or equivalently the ground state of a constrained
      Hamiltonian with a declared boundary at the escape surface;
\item a time window \(t_{\rm micro}\ll t\ll t_{\rm rec}\) on which resonant
      exponential behavior is asserted;
\item a survival amplitude
      \begin{equation}
        A_F(t)=\langle\Psi_F,\ee^{-\ii H_\Lambda t}\Psi_F\rangle
        \label{eq:false-vacuum-survival-amplitude}
      \end{equation}
      or a projected survival probability for the chosen basin;
\item a limiting prescription in volume and regulator, including the order in
      which the semiclassical, long-time, and infinite-volume limits are
      taken.
\end{enumerate}
A decay rate \(\Gamma\) is a statement that, in the declared window,
\[
  A_F(t)=Z_F\ee^{-\ii E_R t-\Gamma t/2}+R_F(t)
\]
with a controlled remainder.  Exact finite-volume spectral decompositions have
real frequencies and long-time recurrences; the exponential law is a resonant
or scaling-window statement.
\end{definition}

This definition is the state-space counterpart of the path-integral contour.
In Euclidean notation the false-vacuum functional is a branch or constrained
functional integral around one basin.  It is not the exact convex Euclidean
1PI effective action of
Chapter~\ref{ch:volume-ii-generating-functionals-1pi}.  The exact convex
Legendre transform mixes competing stable branches when the source is varied.
A bounce calculation instead keeps a chosen false branch and studies its
analytic continuation through the escape saddle.

\paragraph{The negative Gaussian is not an ordinary Gaussian.}
Let \(a\) be the coordinate along a single negative fluctuation direction and
let \(\lambda_-=|\lambda_{\rm neg}|>0\).  Near the saddle the Euclidean weight
has the factor
\[
  \exp\left(+\frac12\lambda_- a^2\right),
\]
which diverges on the real \(a\)-axis.  The steepest-descent cycles are
\(a=\pm \ii y\).  A half-cycle through the saddle gives
\begin{equation}
  \int_{\mathcal J_-^{\rm half}}\dd a\,
  \exp\left(+\frac12\lambda_-a^2\right)
  =
  \pm\frac{\ii}{2}
  \left(\frac{2\pi}{\lambda_-}\right)^{1/2}.
  \label{eq:false-vacuum-negative-mode-half-gaussian}
\end{equation}
The factor of \(1/2\), the sign, and the imaginary unit are contour data.
They are not supplied by the Euclidean equation of motion.  The physical sign
is fixed by the outgoing resonance or by the original \(+\ii0\) deformation of
the false-vacuum energy; choosing the opposite cycle gives the conjugate
resonance.

\begin{example}[Regulated quartic oscillator bounce]
\label{ex:quartic-oscillator-false-vacuum-bounce}
Consider the one-coordinate Euclidean action
\begin{equation}
  S_E[q]
  =
  \int_{-\infty}^{\infty}\dd\tau\,
  \left[
    \frac12\dot q^2
    +
    \frac12\omega^2 q^2
    -
    g q^4
  \right],
  \qquad \omega>0,\quad g>0 .
  \label{eq:quartic-false-vacuum-euclidean-action}
\end{equation}
The potential is unbounded on the real axis, so the false-vacuum problem is
defined by analytic continuation from a stable contour or by a declared pair of
Lefschetz cycles.  The local well at \(q=0\) supplies the false harmonic
state.  Thus this example calibrates the contour and fluctuation normalization
of an analytically continued resonance rather than supplying the
self-adjoint finite-volume datum of
Definition~\ref{def:finite-volume-false-vacuum-datum}.  The Euclidean
equation is
\[
  \ddot q=\omega^2q-4gq^3,
\]
with boundary condition \(q(\tau)\to0\) as
\(|\tau|\to\infty\).  The bounce family is
\begin{equation}
  q_B(\tau-\tau_0)
  =
  \frac{\omega}{\sqrt{2g}}\,
  \operatorname{sech}\!\bigl(\omega(\tau-\tau_0)\bigr).
  \label{eq:quartic-false-vacuum-bounce-profile}
\end{equation}
Using the first integral
\(\frac12\dot q_B^2=V(q_B)\), its action is
\begin{equation}
  B
  =
  S_E[q_B]-S_E[0]
  =
  2\int_0^{\omega/\sqrt{2g}}\dd q\,
    q\sqrt{\omega^2-2gq^2}
  =
  \frac{\omega^3}{3g}.
  \label{eq:quartic-false-vacuum-bounce-action}
\end{equation}
The fluctuation operator is
\begin{equation}
  M_B
  =
  -\frac{\dd^2}{\dd\tau^2}
  +
  \omega^2
  -
  6\omega^2
  \operatorname{sech}^2\!\bigl(\omega(\tau-\tau_0)\bigr).
  \label{eq:quartic-false-vacuum-fluctuation-operator}
\end{equation}
After \(x=\omega(\tau-\tau_0)\), this is the \(l=2\) Pöschl--Teller operator
\(\omega^2[-\partial_x^2+1-6\operatorname{sech}^2x]\).  It has the
node-free negative mode
\[
  \eta_-(x)=\operatorname{sech}^2x,\qquad
  M_B\eta_-=-3\omega^2\eta_-,
\]
and the translational zero mode
\[
  \eta_0(\tau)=\dot q_B(\tau),\qquad M_B\eta_0=0 .
\]
Because the zero mode has one node while the negative mode has none, the
Sturm oscillation theorem gives exactly one negative eigenvalue in this
one-dimensional model.  This proof is special to the regulated scalar
oscillator; in higher-dimensional field theory the corresponding
one-negative-mode statement is an additional spectral assertion about the
chosen bounce and fluctuation operator.

The zero mode is not put into the determinant.  Since
\[
  \int_{-\infty}^{\infty}\dot q_B(\tau)^2\dd\tau=B,
\]
the change from the normalized zero-mode coefficient to the collective
coordinate \(\tau_0\) gives
\begin{equation}
  \frac{\dd c_0}{\sqrt{2\pi}}
  =
  \left(\frac{B}{2\pi}\right)^{1/2}\dd\tau_0 .
  \label{eq:false-vacuum-translation-zero-mode-jacobian}
\end{equation}
Let \(M_F=-\dd_\tau^2+\omega^2\) be the false-harmonic fluctuation operator,
and let \(\det{}''M_B\) denote the determinant with both the negative mode and
the translational zero mode omitted.  In the determinant normalization matched
to the free operator \(M_F\), set
\[
  L_B=-\partial_x^2+1-6\operatorname{sech}^2x .
\]
Its reflectionless scattering phase obeys
\[
  \frac{\dd\delta}{\dd k}
  =
  -2\left(\frac{1}{k^2+1}+\frac{2}{k^2+4}\right),
\]
and the identity
\[
  \int_0^\infty
  \frac{a\log(k^2+1)}{k^2+a^2}\,\dd k
  =
  \pi\log(1+a)
\]
gives the continuum determinant factor
\((2^2 3^2)^{-1}=1/36\).  The dimensionless translation mode
\(\operatorname{sech}x\,\tanh x\) has norm \(2/3\), and the zero-mode
removal formula contributes the inverse square of this norm.  Hence
\[
  |\lambda_-|
  \left(\frac{\det{}''M_B}{\det M_F}\right)_{\rm reg}
  =
  \frac{1}{36}\left(\frac32\right)^2
  =
  \frac{1}{16},
  \qquad |\lambda_-|=3\omega^2 .
\]
Therefore
\begin{equation}
  \left(\frac{\det{}''M_B}{\det M_F}\right)_{\rm reg}
  =
  \frac{1}{48\omega^2}.
  \label{eq:quartic-false-vacuum-positive-determinant-ratio}
\end{equation}
If the absolute value of the negative eigenvalue is included while the zero
mode is still omitted, the corresponding ratio is \(1/16\).  The convention in
\eqref{eq:false-vacuum-one-bounce-amplitude} keeps the negative mode outside
the determinant because its integration cycle supplies the imaginary
half-factor in
\eqref{eq:false-vacuum-negative-mode-half-gaussian}.  A one-bounce
contribution in a time interval of length \(T\) has the form
\begin{equation}
  Z_1(T)
  =
  Z_F(T)\,
  \left(\pm\frac{\ii}{2}\right)
  T
  \left(\frac{B}{2\pi}\right)^{1/2}
  \left(\frac{\det M_F}{\det{}'' M_B}\right)_{\rm reg}^{1/2}
  \ee^{-B}
  \bigl(1+R_1\bigr).
  \label{eq:false-vacuum-one-bounce-amplitude}
\end{equation}
The determinant ratio is a regulated nonzero-mode fluctuation determinant,
not a moduli-space volume.  The residual \(R_1\) contains cubic and higher
fluctuation vertices, determinant-regulator remainders, endpoint errors of the
time interval, and the chosen contour residual.
\end{example}

\paragraph{Dilute exponentiation and real time.}
For \(n\) well-separated bounces the collective coordinates are ordered
\(\tau_1<\cdots<\tau_n\).  If the bounce cores have size \(O(\omega^{-1})\)
and the separation window satisfies
\[
  \omega^{-1}\ll |\tau_i-\tau_j|,\qquad
  n\,\omega^{-1}\ll T,
\]
and if the interaction between separated cores is summably small, the ordered
collective-coordinate volume is \(T^n/n!\) up to a controlled boundary and
interaction residual.  Thus
\begin{equation}
  Z_F^{\rm dilute}(T)
  =
  Z_F(T)
  \exp\left[
    \pm\frac{\ii}{2}
    K_B T\ee^{-B}
  \right]
  \bigl(1+R_{\rm dilute}\bigr),
  \label{eq:false-vacuum-dilute-bounce-exponentiation}
\end{equation}
where
\begin{equation}
  K_B
  =
  \left(\frac{B}{2\pi}\right)^{1/2}
  \left(\frac{\det M_F}{\det{}''M_B}\right)_{\rm reg}^{1/2}
  \label{eq:false-vacuum-bounce-prefactor}
\end{equation}
for the one-coordinate model.  Interpreting
\(Z_F(T)\exp[\pm\ii K_BT\ee^{-B}/2]\) as
\(\exp[-E_{\rm res}T]\) gives
\[
  E_{\rm res}=E_F\mp\frac{\ii}{2}K_B\ee^{-B}.
\]
The outgoing resonance has \(\operatorname{Im}E_{\rm res}<0\), and the
Lorentzian survival probability in the exponential window is
\begin{equation}
  |A_F(t)|^2\simeq \exp[-\Gamma t],
  \qquad
  \Gamma=K_B\ee^{-B}.
  \label{eq:false-vacuum-real-time-width}
\end{equation}
For the analytically continued quartic oscillator above,
\eqref{eq:quartic-false-vacuum-positive-determinant-ratio} gives the leading
resonance width
\begin{equation}
  \Gamma
  =
  \sqrt{\frac{8\omega^5}{\pi g}}\,
  \exp\left[-\frac{\omega^3}{3g}\right]
  \left(1+O\!\left(\frac{g}{\omega^3}\right)\right).
  \label{eq:quartic-false-vacuum-leading-width}
\end{equation}
The passage from the Euclidean branch to a real-time width has two logically
separate state-space inputs.  For a bounded-below finite-volume false vacuum
it uses the metastable state of
Definition~\ref{def:finite-volume-false-vacuum-datum}.  For the unbounded
quartic model it uses the declared outgoing resonance branch supplied by the
analytic continuation or Lefschetz contour.  In either case the calculation
also needs the negative-mode cycle
\eqref{eq:false-vacuum-negative-mode-half-gaussian} and the dilute-gas
residual in \eqref{eq:false-vacuum-dilute-bounce-exponentiation}.

\paragraph{Field-theory extension and proof status.}
For a scalar field in \(D\) Euclidean dimensions, the bounce equations are
\[
  -\partial_E^2\phi_B+V'(\phi_B)=0,\qquad
  \phi_B(x)\to\phi_F\quad |x|\to\infty .
\]
For an \(O(D)\)-symmetric scalar bounce this becomes
\[
  \phi_B''(r)+\frac{D-1}{r}\phi_B'(r)=V'(\phi_B),
  \qquad
  \phi_B'(0)=0,\qquad
  \phi_B(\infty)=\phi_F .
\]
The fluctuation operator is
\begin{equation}
  M_B=-\partial_E^2+V''(\phi_B(x)).
  \label{eq:false-vacuum-field-bounce-fluctuation-operator}
\end{equation}
The \(D\) translational zero modes are separated into the collective coordinate
\(x_0\), with Gram matrix
\[
  G_{\mu\nu}
  =
  \int\dd^Dx\,
  \partial_\mu\phi_B\,\partial_\nu\phi_B .
\]
For an \(O(D)\)-symmetric flat-space bounce, rotational symmetry gives
\[
  G_{\mu\nu}
  =
  \frac{\delta_{\mu\nu}}{D}
  \int\dd^Dx\,(\partial\phi_B)^2 .
\]
With the potential normalized relative to the false vacuum, the scaling
identity for the stationary bounce gives
\[
  \int\dd^Dx\,(\partial\phi_B)^2=DB,
  \qquad
  G_{\mu\nu}=B\,\delta_{\mu\nu}.
\]
The translational collective-coordinate measure therefore contains
\[
  \left(\frac{B}{2\pi}\right)^{D/2}\dd^Dx_0 .
\]
Internal symmetry or gauge zero modes require their own collective
coordinates, stabilizer volumes, gauge fixing, ghost determinants, and orbit
metrics.  The dilute Euclidean bounce gas first supplies an activity per
Euclidean spacetime volume,
\[
  \log Z_E
  \supset
  \pm\frac{\ii}{2}\,\mathcal V_D\,K_{\rm FT}\ee^{-B},
  \qquad
  \mathcal V_D=T V_{D-1}.
\]
After dividing by the observation time \(T\), the physical rate density is
\[
  \frac{\Gamma}{V_{D-1}}
  =
  K_{\rm FT}\ee^{-B}.
\]
Thus \(\mathcal V_D\) is the denominator for Euclidean activity, not the
spatial denominator of the Lorentzian decay rate.  The formula is conditional
on the following data: existence of the bounce in the specified regulator,
exactly one negative mode in the physical scalar sector, separation of all
zero modes, a regulated nonzero-mode determinant ratio, renormalization of
local counterterms, a dilute multi-bounce bound, and a real-time metastable
state or resonance interpretation.  Finite-temperature activation and
gravitational Coleman--De Luccia problems
\cite{ColemanDeLuccia1980} are different problems: the state, boundary
conditions, negative-mode analysis, and contour are not inherited unchanged
from the zero-temperature flat-space calculation.

\section{The Feynman Prescription}

\begin{definition}[Wick-rotated free propagator poles]
\label{def:wick-rotated-free-propagator-poles}
Let
\[
  \omega_{\vec k}=\sqrt{\vec k^2+m^2}.
\]
For a Euclidean separation \((\tau,\vec r)\), with \(\tau>0\), the free
propagator is
\[
  \Delta_E(\tau,\vec r)
  =
  \int\frac{\dd k_\tau\,\dd^{D-1}\vec k}{(2\pi)^D}
  \frac{\ee^{\ii k_\tau\tau+\ii\vec k\cdot\vec r}}
       {k_\tau^2+\omega_{\vec k}^2}.
\]
Continue to complex time by setting \(z^0=\ee^{-\ii\alpha}t\), with
\(0<\alpha\le\pi/2\), and rotate the momentum variable by
\[
  k_\tau=\ii\ee^{\ii\alpha}k^0.
\]
Then
\[
  \Delta_E
  =
  -\ii\ee^{\ii\alpha}
  \int\frac{\dd k^0\,\dd^{D-1}\vec k}{(2\pi)^D}
  \frac{
    \ee^{\ii\vec k\cdot\vec r-\ii k^0t}
  }{
    -\ee^{2\ii\alpha}(k^0)^2+\vec k^2+m^2
  },
\]
with the contour orientation chosen to match the original Euclidean
\(k_\tau\)-contour. The denominator vanishes at
\[
  k^0=\pm \ee^{-\ii\alpha}\omega_{\vec k}.
\]
As \(\alpha\to0^+\), these poles approach the real axis as
\[
  -\omega_{\vec k}+\ii0,
  \qquad
  +\omega_{\vec k}-\ii0.
\]
\end{definition}

Figure~\ref{fig:feynman-prescription-from-wick-rotation} shows how the Euclidean
poles rotate toward the Lorentzian contour and become the Feynman
\(\ii\epsilon\) placement.

\begin{figure}[H]
\centering
\begin{tikzpicture}[>=Stealth,scale=0.95]
  \draw[->,line width=0.7pt] (-3.35,0) -- (3.55,0)
    node[right] {$\operatorname{Re}k^0$};
  \draw[->,line width=0.7pt] (0,-2.25) -- (0,2.35)
    node[above] {$\operatorname{Im}k^0$};
  \draw[purple!70!black,line width=1.1pt,->] (-3.05,0) -- (3.05,0)
    node[below right] {real \(k^0\) contour};
  \node[blue!70!black,scale=1.3] at (0,1.52) {$\times$};
  \node[blue!70!black,right=0.08cm] at (0,1.52) {$-\ee^{-\ii\pi/2}\omega_{\vec k}$};
  \node[blue!70!black,scale=1.3] at (0,-1.52) {$\times$};
  \node[blue!70!black,right=0.08cm] at (0,-1.52) {$+\ee^{-\ii\pi/2}\omega_{\vec k}$};
  \node[violet!70!black,scale=1.25] at (-1.18,1.18) {$\times$};
  \node[violet!70!black,scale=1.25] at (1.18,-1.18) {$\times$};
  \draw[dashed,gray!70,->] (0,1.52) arc[start angle=90,end angle=150,radius=1.52];
  \draw[dashed,gray!70,->] (0,-1.52) arc[start angle=270,end angle=330,radius=1.52];
  \draw[blue!60!black,thick,->]
    (2.05,-0.36) .. controls (1.72,-1.08) and (0.48,-1.36) .. (-0.72,-1.08);
  \draw[blue!60!black,thick,->]
    (-2.05,0.36) .. controls (-1.72,1.08) and (-0.48,1.36) .. (0.72,1.08);
  \node[red!75!black,scale=1.3] at (-2.1,0.17) {$\times$};
  \node[red!75!black,above=0.02cm] at (-2.1,0.17)
    {$-\omega_{\vec k}+\ii0$};
  \node[red!75!black,scale=1.3] at (2.1,-0.17) {$\times$};
  \node[red!75!black,below=0.02cm] at (2.1,-0.17)
    {$+\omega_{\vec k}-\ii0$};
  \node[gray!70!black] at (-2.35,1.10) {$\alpha\to0^+$};
\end{tikzpicture}
\caption{Pole motion under Wick rotation.  The Euclidean poles rotate to the
Lorentzian Feynman positions
\(-\omega_{\vec k}+\ii0\) and \(+\omega_{\vec k}-\ii0\).}
\label{fig:feynman-prescription-from-wick-rotation}
\end{figure}

\paragraph{Feynman boundary value from pole rotation.}
The limiting Lorentzian distribution is therefore
\[
  \Delta_F(x;m^2)
  =
  -\ii
  \int\frac{\dd^Dk}{(2\pi)^D}
  \frac{\ee^{\ii k\cdot x}}{k^2+m^2-\ii\epsilon},
  \qquad
  \epsilon>0,
\]
where
\[
  k\cdot x=-k^0x^0+\vec k\cdot\vec x,
  \qquad
  k^2=-(k^0)^2+\vec k^2.
\]
The limit \(\epsilon\to0^+\) is taken after the distribution has been paired
with test functions or after the contour integral has been evaluated.
Definition~\ref{def:wick-rotated-free-propagator-poles} shows that the
positive-frequency pole approaches the real \(k^0\)-axis from below and the
negative-frequency pole approaches it from above.  The denominator
\(k^2+m^2-\ii\epsilon\), with
\(k^2=-(k^0)^2+\vec k^2\), has precisely these pole placements:
\[
  k^2+m^2-\ii\epsilon
  =
  -(k^0)^2+\omega_{\vec k}^2-\ii\epsilon .
\]
Solving to first order in \(\epsilon\) gives
\[
  k^0=+\omega_{\vec k}-\ii0,
  \qquad
  k^0=-\omega_{\vec k}+\ii0 .
\]
The prefactor \(-\ii\) is the Jacobian and orientation factor obtained from
the rotation \(k_\tau=\ii\ee^{\ii\alpha}k^0\).  The limit is a boundary value
of meromorphic functions in \(k^0\), hence is interpreted only after pairing
with test functions or after an equivalent contour evaluation.

\section{Three Green Functions}

\begin{definition}[Wightman, time-ordered, and Euclidean Green functions]
\label{def:wightman-time-ordered-euclidean-green-functions}
Let \(\widehat\phi\) be a scalar operator-valued distribution with vacuum
\(\vac\). The Wightman \(n\)-point distribution is
\[
  W_n(x_1,\ldots,x_n)
  =
  \bra{\vac}
  \widehat\phi(x_1)\cdots\widehat\phi(x_n)
  \ket{\vac}.
\]
The order displayed in the formula is part of the definition.

For distinct time coordinates, the time-ordered Green function is
\[
  G_{T,n}(x_1,\ldots,x_n)
  =
  \bra{\vac}
  \operatorname{T}\{
    \widehat\phi(x_1)\cdots\widehat\phi(x_n)
  \}
  \ket{\vac}.
\]
If
\[
  x_{i_1}^0>x_{i_2}^0>\cdots>x_{i_n}^0,
\]
then
\[
  G_{T,n}(x_1,\ldots,x_n)
  =
  W_n(x_{i_1},\ldots,x_{i_n}).
\]
Coincident-time prescriptions are contact-term data and must be fixed in the
chosen operator or path-integral framework.

The Euclidean Green function is a distribution in Euclidean points
\(x_{E,a}=(\tau_a,\vec x_a)\).  It is denoted \(G_{E,n}\) in this
path-integral construction.  When the same distributions are used as OS
reconstruction data, they are denoted \(S_n\) and called Schwinger functions.
When the analytic continuation from Wightman functions exists and
\[
  \tau_{i_1}>\tau_{i_2}>\cdots>\tau_{i_n},
\]
it is given by
\[
  G_{E,n}(x_{E,1},\ldots,x_{E,n})
  =
  W_n(x_{i_1},\ldots,x_{i_n})
  \big|_{x_a^0=-\ii\tau_a}.
\]
\end{definition}

These three distributions encode different boundary values of the same
analytic structure when the required analyticity and spectral assumptions are
available. The next chapter uses this structure in the two-point function to
derive the K\"all\'en--Lehmann spectral representation.
